The command-line menus remain fully supported and are the
authoritative feature inventory, but they are frozen.  New interactive
work normally uses the GUI described in Chapter
\ref{modern_architecture}; new programmatic work uses the REST API in
Chapter \ref{rest_api}.

\section{Reproducible sessions}

Automated runs write every menu answer to a text file, one answer per
line, and feed that file through standard input.  Always supply a seed
and capture standard output:

\begin{verbatim}
./build/neuron --seed 42 < session.in > session.out
\end{verbatim}

Read the completed transcript and verify that it ends after the quit
confirmation, not at an error or unanswered prompt.  Do not automate
the menus by typing into a live terminal.

\section{Menu organization}

\begin{description}
\item[Specify dataset] Set input and output counts; load raw data or a
normalized training and test pair; create a seeded train and test
split; set outcome type, normalization bounds, and ROC reporting; save
normalized sets and scaling factors.
\item[Specify model] Create a feed-forward network or binary logistic
model; set bias nodes, hidden layers, and error function; load a saved
network; set logging.
\item[Use model] Randomize weights; configure learning rate, weight
decay, batch or epoch behavior, printing, and stopping conditions;
train with canonical gradient descent, conjugate gradient descent, or
Shanno's method; run stepwise selection; save the network and
predictions.
\item[Discriminant function analysis] Run linear or quadratic
discriminant analysis and save its predictions.
\end{description}

The exact menu-to-GUI-to-API mapping is maintained in
\path{docs/gui_cli_parity.md}.  That file is the current
option-level reference and is updated whenever an interface changes.

\section{Dataset and run files}

Numeric data place the outcome in the final column.  Binary outcomes
are encoded as 0 and 1.  The first loaded dataset establishes the run
directory, where the engine writes networks, scaling factors, reports,
predictions, \texttt{neuron.log}, and
\texttt{neuron\_actions.log}.

For raw CSV data, use \texttt{tools/mkdataset.py} as described in
Chapter \ref{utility_chap}.  When an intercept-bearing model will be
trained, reference-category coding avoids collinearity between a full
one-hot group and the intercept.

\section{Convergence and reporting}

The iteration ceiling is a safety limit, not evidence of convergence.
If training reaches it, the current weights can be continued, but the
fit is unfinished.  Model-selection and cross-validation procedures
refuse unfinished candidate fits.

For binary outcomes, interpret held-out performance through
classification measures and ROC areas, not accuracy alone.  The
primary fitted binormal $A_z$ is reported with its bootstrap 95\%
confidence interval, operating-point count, resample count, and any
failures.  The exact empirical AUC and Hanley--McNeil interval are an
independent cross-check.  Logistic reports additionally provide Wald
tests and a condition number.

\section{Current command-line options}

\begin{verbatim}
Usage: neuron [--seed N] [--gui [--no-browser]] [--version]
\end{verbatim}

\texttt{--gui} starts the local browser interface.
\texttt{--no-browser} leaves it running without opening a browser.
\texttt{--version} prints the package version and exits.
